\begin{multicols}{2}
~\\
On considère le prisme droit ABFEDCGH tel que $\mathrm{AB} = \mathrm{AD}$.\\

Sa base ABFE est un trapèze rectangle en A, vérifiant $\vect{\mathrm{BF}} = \dfrac{1}{2} \vect{\mathrm{AE}}$.\\

On note I le milieu du segment [EF].\\

On note J le milieu du segment [AE].\\

On associe à ce prisme le repère orthonormé $\left(\text{A}~;~ \vect{\imath\vphantom{k}}, \vect{\jmath\vphantom{k}}, \vect{k}\right)$ tel que :\\
$\vect{\imath} = \vect{\mathrm{AB}} ;\quad \vect{\jmath} = \vect{\mathrm{AD}} ;\quad \vect{k} = \vect{\mathrm{AJ}}$

\vfill\null

\columnbreak

\scalebox{0.7}{
\psset{unit=1cm,radius=2pt}
\def\xmin{-1} \def\xmax{8} \def\ymin{-1} \def\ymax{10.5}
\begin{pspicture*}(\xmin,\ymin)(\xmax,\ymax)

\Cnode*(0,0){A} \Cnode*(4,-0.5){B} \Cnode*(7,0.5){C} \Cnode*(3,1){D} 
\uput[dl](A){A} \uput[d](B){B} \uput[r](C){C} \uput[ur](D){D} 
\Cnode*(0,9){E} \Cnode*(4,4){F} \Cnode*(7,5){G} \Cnode*(3,10){H}
\uput[ul](E){E} \uput[80](F){F} \uput[60](G){G} \uput[u](H){H} 
\Cnode*[radius=2pt](0,4.5){J} \Cnode*[radius=2pt](2,6.5){I} 
\uput[r](J){J} \uput[ur](I){I}
\psline(B)(C)(G)(F)(B)(A)(E)(H)(G) \psline(E)(F)
\psline[linestyle=dashed](C)(D)(H) \psline[linestyle=dashed](A)(D)
\end{pspicture*}
}
\end{multicols}

\begin{enumerate}

\item On donne les coordonnées de quatre vecteurs dans la base $\left(\vect{\imath\vphantom{k}}, \vect{\jmath\vphantom{k}}, \vect{k}\right)$. \\
Lequel est un vecteur normal au plan (ABG)?

	\begin{tabularx}{\linewidth}{*{4}{>{\bfseries\arraybackslash}p{5mm}@{}X}}
		a.&$\vect{n} \begin{pmatrix}1 \\ 1 \\ 1\end{pmatrix}$&
		b.&$\vect{n} \begin{pmatrix}-1 \\ 1 \\ 1\end{pmatrix}$&
		c.&$\vect{n} \begin{pmatrix}0 \\ -1 \\ 1\end{pmatrix}$&
		d.&$\vect{n} \begin{pmatrix}0 \\ 0 \\ 1\end{pmatrix}$\\
	\end{tabularx}

Le plan (ABG) a pour vecteurs directeurs $\vectt{AB}$ et $\vectt{BG}$.

$\vectt{AB}$ a pour coordonnées $\begin{pmatrix}1 \\ 0 \\ 0\end{pmatrix}$ et\\
$\vectt{BG} = \vectt{BC}+\vectt{CG} = \vectt{AD}+\vectt{AJ}$ donc $\vectt{BG}$ a pour coordonnées  $ \begin{pmatrix} 0 \\ 1 \\ 1 \end{pmatrix}$.

On cherche donc, parmi les trois vecteurs proposés, celui qui est orthogonal à la fois à $\vectt{AB}$ et à $\vectt{BG}$.
Pour $\vect{n} \begin{pmatrix}0 \\ -1 \\ 1\end{pmatrix}$, on a:\\
$\vect{n}\cdot\vectt{AB}= 0\times 1+(-1)\times 0+1\times 0=0$ donc $\vect{n}\perp \vectt{AB}$;\\
$\vect{n}\cdot\vectt{BG}= 0\times 0+(-1)\times 1+1\times 1=0$ donc $\vect{n}\perp \vectt{BG}$

\hfill\textbf{Réponse c.}

	\item Parmi les droites suivantes, laquelle est parallèle à la droite (IJ) ?

	\begin{tabularx}{\linewidth}{*{4}{>{\bfseries\arraybackslash}p{5mm}@{}X}}
		a.&(DG)&
		b.&(BD)&
		c.&(AG)&
		d.&(FG)\\
	\end{tabularx}

La droite (IJ) est parallèle à la droite (AF) (droite des milieux dans le triangle AEF).

Les vecteurs $\vectt{AF}$ et $\vectt{DG}$ sont égaux donc les droites (AF) et (DG) sont parallèles.

On en déduit que  les droites (IJ) et (DG) sont parallèles.

\hfill\textbf{Réponse a.}

	\item Quels vecteurs forment une base de l'espace ?

	\begin{tabularx}{\linewidth}{*{4}{>{\bfseries\arraybackslash}p{5mm}@{}X}}
		a.& $\left(\vect{\mathrm{AB}} ; \vect{\mathrm{CG}}\right)$&
		b.& $\left(\vect{\mathrm{AB}} ; \vect{\mathrm{AC}} ; \vect{\mathrm{AD}}\right)$ &
		c.& $\left(\vect{\mathrm{DA}} ; \vect{\mathrm{DC}} ; \vect{\mathrm{DG}}\right)$&
		d.& $\left(\vect{\mathrm{CA}} ; \vect{\mathrm{CG}} ; \vect{\mathrm{CE}}\right)$\\
	\end{tabularx}

Il faut trouver trois vecteurs non coplanaires.

On peut éliminer la proposition \textbf{a} car il n'y a que deux vecteurs.

On peut éliminer la proposition \textbf{b} car $\vectt{AC} = \vectt{AB} + \vectt{AD}$ donc les trois vecteurs sont coplanaires. 

On peut éliminer la proposition \textbf{d} car $\vectt{CE} = \vectt{CA} + \vectt{AE} = \vectt{CA} + 2\vectt{CG}$ donc les trois vecteurs sont coplanaires. 

\hfill\textbf{Réponse c.}

	\item Une décomposition du vecteur $\vect{\mathrm{AG}}$ comme somme de plusieurs vecteurs \textbf{deux à deux orthogonaux} est :

	\begin{tabularx}{\linewidth}{*{2}{>{\bfseries\arraybackslash}p{5mm}@{}X}}
		a.& $\vect{\mathrm{AG}} = \vect{\mathrm{AB}} + \vect{\mathrm{HG}}$ &
		b.& $\vect{\mathrm{AG}} = \vect{\mathrm{AB}} + \vect{\mathrm{AD}} + \vect{\mathrm{AJ}}$\\
		c.& $\vect{\mathrm{AG}} = \vect{\mathrm{AB}} + \vect{\mathrm{BJ}} + \vect{\mathrm{JG}}$&
		d.& $\vect{\mathrm{AG}} = \vect{\mathrm{AD}} + \vect{\mathrm{DH}} + \vect{\mathrm{HG}}$\\
	\end{tabularx}

$\vectt{AG} = \vectt{AB} +  \vectt{BG}$ et $ \vectt{BG} \neq  \vectt{HG}$; donc la proposition \textbf{a} est fausse.

Les trois vecteurs $ \vectt{AB}$,  $\vectt{AD}$ et  $\vectt{AJ}$ sont orthogonaux deux à deux puisqu'ils forment une base orthonormée.

De plus $ \vectt{AB}+\vectt{AD}+\vectt{AJ} = \vectt{AB}+\vectt{BC}+\vectt{CG}= \vectt{AG}$.

\hfill\textbf{Réponse b.} 

	\item Le volume du prisme droit ABFEDCGH, est égal à :

		\begin{tabularx}{\linewidth}{*{4}{>{\bfseries\arraybackslash}p{5mm}@{}X}}
			a.& $\dfrac{5}{8}$&
			b.& $\dfrac{8}{5}$&
			c.& $\dfrac{3}{2}$&
			d.& 2\\
		\end{tabularx}

Si on appelle K le milieu de [DH], on peut dire que le  prisme est composé du cube ABCDJFGK, de volume 1, surmonté du prisme JFGKHE, qui est la moitié du cube du dessous, donc qui a pour volume $\dfrac{1}{2}$. Le volume total fait donc $\dfrac{3}{2}$.

\hfill\textbf{Réponse c.}
\end{enumerate}
